\documentclass[11pt,a4paper]{article}
\usepackage[utf8]{inputenc}
\usepackage[T1]{fontenc}
\usepackage[italian]{babel}
\usepackage[margin=2.5cm]{geometry}
\usepackage{xcolor}
\usepackage{enumitem}
\usepackage{hyperref}
\usepackage{fancyhdr}
\usepackage{titling}
\usepackage{tabularx}
\usepackage{booktabs}
\usepackage{textcomp}

\definecolor{studioteal}{HTML}{0D6E6E}
\definecolor{studiogrigio}{HTML}{5B6B6A}
\definecolor{studioarancio}{HTML}{B45309}

\hypersetup{
  colorlinks=true,
  linkcolor=studioteal,
  urlcolor=studioteal,
  citecolor=studioteal,
  pdftitle={Guida per lo staff dello studio},
  pdfauthor={Studio Medico Dottoressa Braglia}
}

\pagestyle{fancy}
\fancyhf{}
\fancyhead[L]{\small\color{studiogrigio} Studio Medico Dottoressa Braglia --- Pannello}
\fancyhead[R]{\small\color{studiogrigio} Guida per lo staff}
\fancyfoot[C]{\small\thepage}
\renewcommand{\headrulewidth}{0.4pt}
\renewcommand{\headrule}{\hbox to\headwidth{\color{studioteal}\leaders\hrule height \headrulewidth\hfill}}

\setlist[itemize]{leftmargin=1.4em, itemsep=0.25em}
\setlist[enumerate]{leftmargin=1.6em, itemsep=0.35em}

\newcommand{\schermata}[1]{\textbf{\color{studioteal}#1}}
\newcommand{\bottone}[1]{\textbf{``#1''}}
\newcommand{\solomedico}{\textcolor{studioarancio}{\textbf{[solo medico]}}\ }

\title{\color{studioteal}\Huge Guida per lo staff}
\author{\Large Studio Medico Dottoressa Braglia}
\date{}

\begin{document}

\begin{titlepage}
\centering
\vspace*{3cm}
{\color{studioteal}\Huge\bfseries Guida per lo staff}\\[1em]
{\Large Come usare il pannello di gestione}\\[2em]
{\large Studio Medico Dottoressa Braglia}\\
{\large Ambulatori di Arceto e Casalgrande}\\[3em]
{\normalsize\color{studiogrigio} \url{https://studiomedicobragliavaleria.it/admin.html}}
\vfill
{\small\color{studiogrigio}
Questa guida è per chi lavora nello studio: medico e segreteria.\\
Le voci segnate \solomedico sono visibili solo all'accesso del medico.}
\end{titlepage}

\tableofcontents
\newpage

\section{Primo accesso}

Il pannello è raggiungibile all'indirizzo \url{https://studiomedicobragliavaleria.it/admin.html}. È separato dal sito che vedono i pazienti: ha un proprio indirizzo e un proprio accesso.

Un nuovo collaboratore riceve una \textbf{password provvisoria}: compare a schermo quando il medico crea l'accesso, e arriva anche via email insieme al link del pannello e a questa stessa guida in allegato. Al primo ingresso il sistema chiede subito di sceglierne una personale, che nessun altro conoscerà: finché non viene scelta, il resto del pannello resta bloccato.

\section{Le schede del pannello}

In alto si trovano le schede principali. Alcune sono visibili solo al medico.

\begin{center}
\begin{tabularx}{\linewidth}{lX}
\toprule
\textbf{Scheda} & \textbf{A cosa serve} \\
\midrule
Oggi & Riepilogo della giornata: quante visite, cosa aspetta una risposta. \\
Prenotazioni & Agenda, richieste da confermare, prenotazioni a mano. \\
Calendario & Vista mensile: un colpo d'occhio su quali giorni sono più pieni. \\
Chiusure & Bloccare un giorno o una fascia oraria alle prenotazioni. \\
Farmaci & Richieste di rinnovo ricette. \\
Visite specialistiche & Richieste di impegnativa per uno specialista. \\
Esami del sangue & Richieste di impegnativa per le analisi. \\
Pazienti & Anagrafica, ricerca, fascicolo, dati e password di accesso. \\
Collaboratori & \solomedico Chi ha un accesso al pannello, e con quale ruolo. \\
Impostazioni & Cambiare la propria password quando si vuole. \\
Stato del sistema & Email, backup, coda degli invii. \\
\bottomrule
\end{tabularx}
\end{center}

\section{Oggi}

La prima schermata dopo l'accesso. Mostra a colpo d'occhio, in una serie di riquadri cliccabili: quante visite ci sono oggi, quante richieste di visita dal sito sono ancora da confermare, quante visite future sono già confermate, quante richieste di farmaci/visite specialistiche/esami aspettano una risposta, e quanti pazienti sono in archivio. Ogni riquadro porta dritto alla scheda corrispondente, già filtrata, e sotto compare anche l'agenda della giornata.

\section{Calendario}

Una vista d'insieme del mese, utile per farsi un'idea di quali giorni sono più pieni prima ancora di aprire l'agenda giorno per giorno: ogni giorno con almeno una prenotazione porta un pallino (tenendoci sopra il puntatore si vede quante sono). Cliccando su un giorno si apre il dettaglio con orario, paziente, telefono, ambulatorio, motivo, stato e codice di ogni visita.

\section{Prenotazioni}

\subsection{Confermare o rifiutare una richiesta dal sito}

Le visite prenotate dal sito o dal chatbot dei pazienti nascono come \emph{richieste}, in attesa: appaiono con il filtro \bottone{Da confermare}. Aprendo una richiesta si può:

\begin{itemize}
  \item \textbf{Confermarla} così com'è (o cambiando giorno, ora o ambulatorio se necessario) --- al paziente arriva l'email di conferma definitiva;
  \item \textbf{Rifiutarla}, indicando il motivo --- il paziente lo legge nell'email di risposta, e il posto torna libero.
\end{itemize}

\subsection{Prenotare a mano (telefono, sportello)}

Dal pulsante \bottone{+ Nuova prenotazione} si registra una visita presa al telefono o allo sportello. Mentre si scrive nome, cognome, telefono o email, il pannello suggerisce i pazienti già in archivio che corrispondono: scegliendone uno, gli altri campi si riempiono da soli e la visita resta agganciata a quella scheda, senza creare un doppione.

La casella \bottone{orario mio / forza} permette di scavalcare i limiti pensati per chi prenota da solo (orari di apertura, quanti giorni in anticipo, il passato): utile per registrare una visita già avvenuta, o un caso particolare. Non si può mai, nemmeno forzando, mettere due pazienti nello stesso posto alla stessa ora.

\subsection{Spostare o annullare}

Una prenotazione confermata si può riprogrammare (cambiare giorno, ora o ambulatorio) o annullare. In entrambi i casi parte un'email al paziente, e chi era in lista d'attesa per quel giorno viene avvisato se si libera un posto.

Quando si sposta un appuntamento si può aggiungere un \bottone{Motivo dello spostamento} (facoltativo): se lo si scrive, il paziente lo legge nella sua email di conferma --- utile quando il cambio è dovuto a un'esigenza dello studio (``il medico non è disponibile quel giorno'') e non a un semplice orario più comodo, che il paziente non si aspetta di dover chiedere per telefono.

\subsection*{Dove restano visibili gli annullamenti e i rifiuti}

Una prenotazione appena annullata o rifiutata non sparisce subito dall'elenco: resta in cima, ben visibile, per le \textbf{24 ore} successive --- utile se un paziente richiama subito dopo chiedendo conferma dell'annullamento. Le richieste ancora \bottone{Da confermare} restano comunque le prime in assoluto, sopra agli annullamenti recenti: sono quelle che aspettano un'azione. Lo stesso vale per le code di farmaci, visite specialistiche ed esami.

\section{Chiusure}

Per bloccare le prenotazioni --- ferie, un imprevisto, mezza giornata --- si aggiunge una chiusura: un giorno intero, un intervallo di giorni, oppure solo una fascia oraria all'interno di una o più giornate. Le prenotazioni già confermate in quel periodo \textbf{non} vengono annullate da sole: il pannello elenca quelle da avvisare a mano.

\section{Farmaci, visite specialistiche, esami del sangue}

Le tre code funzionano allo stesso modo. Per ogni richiesta nuova si può:

\begin{itemize}
  \item \textbf{Confermarla}, indicando se serve il numero della ricetta o dell'impegnativa --- il paziente riceve l'email con le istruzioni per il ritiro;
  \item \textbf{Modificarla}, se dopo una telefonata è cambiato qualcosa rispetto a quanto chiesto --- l'email al paziente mostra cosa aveva chiesto e cosa gli è stato preparato;
  \item \textbf{Rifiutarla}, con un motivo.
\end{itemize}

Se il paziente ha allegato la foto di una prescrizione, è visibile direttamente nella scheda della richiesta.

Anche qui, dal pulsante \bottone{+ Aggiungi}, si registra una richiesta presa al telefono, con lo stesso suggerimento dei pazienti già in archivio.

\section{Pazienti}

\subsection{Elenco e ricerca}

Mostra tutti i pazienti in carico allo studio. Chi non è più assistito non compare, a meno di spuntare l'opzione apposita: dimettere qualcuno serve proprio a toglierlo dagli elenchi di tutti i giorni, senza perdere la sua storia.

La ricerca funziona scrivendo più parole insieme, in qualunque ordine e con qualunque maiuscola/minuscola: ``mario rossi'', ``Rossi Mario'' o ``ROSSI mario'' trovano tutti la stessa persona, anche se nome e cognome sono scritti in campi separati.

\subsection{Aggiungere un paziente a mano}

Il pulsante \bottone{+ Aggiungi paziente}, in cima all'elenco, serve per inserire una persona che non è mai passata dal sito. Bastano nome e cognome; telefono ed email sono facoltativi.

Dalla scheda \bottone{Dati e password} di un paziente già in archivio c'è anche \bottone{+ Aggiungi un familiare}: pensato per chi --- una persona anziana, per esempio --- non ha un proprio telefono o una propria email. Il modulo propone di \textbf{riusare gli stessi contatti} del paziente già aperto (si può togliere la spunta e scriverne di diversi).

Se alla fine c'è un'email in archivio (propria o condivisa), \textbf{l'accesso al sito si apre nello stesso passaggio}: compare a schermo la password provvisoria, da copiare o annotare subito --- non sarà più recuperabile in seguito. Se l'email è condivisa con un altro accesso già esistente, il pannello lo segnala: restano due account distinti, separati solo dalla password con cui si accede (il paziente, dal suo \bottone{Modifica account}, trova poi il familiare in \bottone{Gestisci come} e può passare dall'uno all'altro senza reinserire nessuna password).

\subsection{Il fascicolo} \solomedico

Cliccando sul nome si apre il fascicolo completo: le visite passate e future, i medicinali che il paziente prende abitualmente, e la storia di ogni richiesta con come è stata gestita. Contiene il motivo delle visite, un dato clinico: per questo lo apre solo il medico.

\subsection{Dati e password}

Sotto ogni nome, il pulsante \bottone{Dati e password} --- questo lo usa chiunque lavori nello studio, medico o segreteria. Da lì si può:

\begin{itemize}
  \item correggere nome, cognome, telefono o email del paziente;
  \item \bottone{Reimposta password} --- se il paziente ha già un accesso e non riesce più a entrare; oppure \bottone{Crea accesso al sito} al suo posto, se non ne ha ancora uno e c'è un'email in archivio;
  \item \bottone{Riattiva email su Brevo} --- per chi ha cliccato per errore ``annulla iscrizione'' su un'email automatica: da quel momento non riceve più nemmeno le conferme delle prenotazioni, finché non lo si riattiva da qui;
  \item \bottone{+ Aggiungi un familiare} --- descritto sopra.
\end{itemize}

Dopo ogni correzione o nuovo accesso, il paziente riceve un'email che lo avvisa. Se il paziente non ha ancora un'email in archivio, non si può ancora aprirgli un accesso: il pannello lo segnala chiaramente, e basta scriverla prima nei dati.

\subsection{Dimettere o cancellare}

Due gesti diversi, apposta:

\begin{itemize}
  \item \bottone{Non è più nostro paziente} --- sparisce dagli elenchi, la sua storia resta in archivio. Si annulla con un altro clic (``È tornato nostro paziente''). È il gesto giusto quasi sempre.
  \item \bottone{Cancella definitivamente} \solomedico --- non si torna indietro: visite, richieste, ricette e foto allegate spariscono per sempre, insieme all'eventuale accesso del paziente al sito. Va usata solo per un doppione o su richiesta esplicita della persona. Il pannello mostra prima quanto andrebbe perso, e chiede di scrivere il cognome per confermare.
\end{itemize}

\section{Collaboratori} \solomedico

Solo il medico può creare, sospendere o eliminare gli accessi al pannello.

\begin{itemize}
  \item \textbf{Creare un accesso}: nome, email, e il ruolo (\emph{segreteria} o \emph{medico}). Viene generata una password provvisoria, da consegnare al collaboratore: compare una volta sola sullo schermo, quindi va copiata o annotata subito.
  \item \textbf{Sospendere / riattivare}: chiude fuori subito un accesso, anche se la persona è già collegata.
  \item \textbf{Rinnovare la password}: genera una nuova password provvisoria per chi l'ha persa, e chiude le sessioni aperte con quella vecchia.
  \item \textbf{Eliminare}: toglie l'accesso definitivamente.
\end{itemize}

Non è possibile eliminare, sospendere o togliere i poteri al proprio stesso accesso, né lasciare lo studio senza nemmeno un accesso da medico attivo: il sistema lo impedisce apposta, per non restare chiusi fuori dal pannello.

\section{Impostazioni}

Ogni collaboratore --- medico o segreteria --- può cambiare la propria password quando vuole da questa scheda, non solo la prima volta con quella provvisoria.

\section{Notifiche e installazione come app}

In alto, accanto al pulsante ``Esci'', sono disponibili:

\begin{itemize}
  \item \bottone{Attiva notifiche} --- per ricevere un avviso sul dispositivo quando arriva una nuova richiesta da confermare (visita, farmaci, specialistica, esami). Il testo del pulsante cambia da solo: \bottone{Disattiva notifiche} quando sono accese, \bottone{Notifiche disattivate} se sono state bloccate dalle impostazioni del telefono (si riattivano da lì, non dal pannello). Quando sono accese compare anche la casella \bottone{Suono}, per ricevere gli avvisi in silenzio su quel dispositivo senza spegnerli del tutto.
  \item \bottone{Installa} --- per aggiungere un'icona del pannello sul telefono o computer, che si apre a schermo intero come un'app vera.
\end{itemize}

Su iPhone e iPad, per entrambe le funzioni, va prima aggiunto il pannello alla schermata Home da Safari (icona Condividi \textrightarrow\ ``Aggiungi alla schermata Home''): è una regola di Apple, non una nostra scelta, e vale anche usando Chrome.

\section{L'assistente del pannello}

In basso a destra è disponibile un assistente che risponde a domande scritte in linguaggio naturale, sempre con numeri veri presi dal database (mai inventati). Alcuni esempi:

\begin{center}
\begin{tabularx}{\linewidth}{lX}
\toprule
\textbf{Si scrive} & \textbf{Cosa fa} \\
\midrule
``come va oggi'' & Il riepilogo della giornata. \\
``cosa devo vedere'' & Solo ciò che aspetta ancora una risposta. \\
``chi viene oggi'' / ``chi viene domani'' & L'agenda del giorno. \\
``quante prenotazioni ho oggi'' & Un numero solo, súbito. \\
``cerca Rossi'' & Trova la persona e apre la scheda Pazienti. \\
\emph{(incollare un codice, es. PRE-A1B2-C3D4)} & Dice cos'è e a che punto sta. \\
``annulla PRE-1234-ABCD'' & Annulla l'appuntamento (chiede conferma). \\
``blocca il 15/10 dalle 10 alle 12'' & Crea una chiusura (chiede conferma). \\
``che chiusure ci sono'' & L'elenco delle chiusure future. \\
``apri i farmaci'' & Porta dritto sulla scheda giusta. \\
``password'', ``notifiche'', ``installa'' & Spiega dove trovarle. \\
\bottomrule
\end{tabularx}
\end{center}

Il pulsante \bottone{Al telefono} apre lo stesso modulo che vede il paziente, per registrare una richiesta raccolta a voce.

\section{Stato del sistema}

Mostra se l'invio delle email è attivo, quante ce ne sono ancora in coda (le email non si perdono mai: se l'invio non è disponibile, restano lì e riprovano da sole), e lo stato dell'ultima copia di sicurezza del database.

\section{Promemoria automatici}

Il giorno prima di ogni visita confermata, il paziente riceve automaticamente un promemoria via email, con la richiesta di annullare se non può presentarsi.

\section{Buone pratiche}

\begin{itemize}
  \item Non condividere mai la propria password con altri: ogni accesso è personale, e serve a sapere chi ha fatto cosa.
  \item Sospendere subito l'accesso di chi lascia lo studio.
  \item Usare \bottone{Cancella definitivamente} solo quando è davvero irreversibile: nel dubbio, \bottone{Non è più nostro paziente} si annulla, l'altra no.
  \item Per un cambio di email o telefono di un paziente, preferire ``Dati e password'' invece di far ricreare un account nuovo: mantiene collegata tutta la sua storia.
\end{itemize}

\end{document}
